\chapter*{References}
\addcontentsline{toc}{chapter}{References}

\refitem{Ai, X., Wang, Y., Wang, P., \& Wang, S. (2025). Impact of visual distractors in virtual reality environments on sustained attention behavioral performance and EEG characteristics. \textit{Frontiers in Human Neuroscience}. \url{https://pmc.ncbi.nlm.nih.gov/articles/PMC12698649/}}

\refitem{Apple (2024). \textit{Accessing the main camera — visionOS enterprise APIs.} \url{developer.apple.com/documentation/visionOS/accessing-the-main-camera}}

\refitem{Bailey, R., McNamara, A., Sudarsanam, N., \& Grimm, C. (2009). Subtle gaze direction. \textit{ACM Transactions on Graphics, 28}(4). \url{https://doi.org/10.1145/1559755.1559757}}

\refitem{Ball, K., \& Owsley, C. (1993). The useful field of view test: a new technique for evaluating age-related declines in visual function. \textit{Journal of the American Optometric Association, 64}(1), 71–79. PMID: 8454831.}

\refitem{Barhorst-Cates, E. M., Rand, K. M., \& Creem-Regehr, S. H. (2016). The Effects of Restricted Peripheral Field-of-View on Spatial Learning while Navigating. \textit{PLoS ONE}, 11(10), e0163785.}

\refitem{Biocca, F., Tang, A., Owen, C., \& Xiao, F. (2006). Attention funnel: omnidirectional 3D cursor for mobile augmented reality platforms. \textit{Proc. CHI 2006}. \url{https://ieeexplore.ieee.org/document/1579336/}}

\refitem{Brown, P., Spronck, P., \& Powell, W. (2022). The simulator sickness questionnaire, and the erroneous zero baseline assumption. \textit{Frontiers in Virtual Reality, 3}:945800. \url{https://doi.org/10.3389/frvir.2022.945800}}

\refitem{Buchner, J., Buntins, K., \& Kerres, M. (2022). The Impact of Augmented Reality on Cognitive Load and Performance: A Systematic Review. \textit{Journal of Computer Assisted Learning}, 38(1), 285–303. doi:10.1111/jcal.12617}

\refitem{Caine, K. (2016). Local standards for sample size at CHI. In \textit{Proceedings of CHI 2016} (pp. 981–992). ACM. \url{https://dl.acm.org/doi/10.1145/2858036.2858498}}

\refitem{Cao, Z., Grandi, J., \& Kopper, R. (2021). Granulated Rest Frames Outperform Field of View Restrictors on Visual Search Performance. \textit{Frontiers in Virtual Reality}, 2:604889. doi:10.3389/frvir.2021.604889}

\refitem{Cheng, Y., Yin, H., Yan, Y., Gugenheimer, J., \& Lindlbauer, D. (2022). Towards understanding diminished reality. In \textit{Proceedings of CHI 2022}. ACM. \url{https://dl.acm.org/doi/10.1145/3491102.3517452}}

\refitem{Coviello, R. [xrdevrob] (2025). \textit{QuestCameraKit} [software]. GitHub. \url{https://github.com/xrdevrob/QuestCameraKit}}

\refitem{Duan, H., et al. (2022). Saliency in Augmented Reality (SARD dataset). \textit{ACM Multimedia 2022}.}

\refitem{Erickson, A., Bruder, G., \& Welch, G. (2023). Analysis of the Saliency of Color-Based Dichoptic Cues in Optical See-Through Augmented Reality. \textit{IEEE Transactions on Visualization and Computer Graphics.} doi:10.1109/TVCG.2022.3195111}

\refitem{Exploring Diminished Reality for Attention Support: A Co-Design Study with Students with ADHD. \textit{DIS 2026}. \url{https://dl.acm.org/doi/10.1145/3800645.3813095}}

\refitem{Fernandes, A. S., \& Feiner, S. K. (2016). Combating VR Sickness through Subtle Dynamic Field-of-View Modification. \textit{IEEE 3DUI 2016}, 201–210. doi:10.1109/3DUI.2016.7460053}

\refitem{Grogorick, S., Stengel, M., Eisemann, E., \& Magnor, M. (2017). Subtle Gaze Guidance for Immersive Environments. \textit{ACM SAP 2017}. doi:10.1145/3119881.3119890}

\refitem{Grogorick, S., Albuquerque, G., \& Magnor, M. (2018). Comparing Unobtrusive Gaze Guiding Stimuli in Head-Mounted Displays. \textit{IEEE ICIP 2018}, 2805–2809. doi:10.1109/ICIP.2018.8451784. Companion: \textit{ACM TAP}, 15(4), 2018. doi:10.1145/3238303}

\refitem{Hart, S. G., \& Staveland, L. E. (1988). Development of NASA-TLX. In \textit{Human Mental Workload}.}

\refitem{Hata, H., Koike, H., \& Sato, Y. (2016). Visual Guidance with Unnoticed Blur Effect. \textit{AVI 2016}, 28–35. doi:10.1145/2909132.2909254}

\refitem{Hein, P., Bernhagen, M., \& Bullinger, A. C. (2019). Improving Visual Attention Guiding by Differentiation between Fine and Coarse Navigation. \textit{2019 11th International Conference on Virtual Worlds and Games for Serious Applications (VS-Games)}. doi:10.1109/VS-Games.2019.8864539}

\refitem{Higgins, P., Barron, R., \& Matuszek, C. (2022). Head Pose as a Proxy for Gaze in Virtual Reality. \textit{VAM-HRI 2022}. \url{https://iral.cs.umbc.edu/Pubs/Higgins2022VAM-HRI.pdf}}

\refitem{Hong, J., Langlotz, T., Sutton, J., \& Regenbrecht, H. (2024). Visual Noise Cancellation: Exploring Visual Discomfort and Opportunities for Vision Augmentations. \textit{ACM Transactions on Computer-Human Interaction.} doi:10.1145/3634699}

\refitem{Immersed (2022). \textit{Passthrough Portals.} \url{uploadvr.com/immersed-adds-passthrough-portals-tracked-keyboards/}}

\refitem{ISO (2016). \textit{ISO 17488:2016 — Road vehicles — Transport information and control systems — Detection-response task (DRT) for assessing attentional effects of cognitive load in driving.}}

\refitem{ISO 17488:2016. Road Vehicles — Detection-Response Task (DRT) for Assessing Attentional Effects of Cognitive Load in Driving.}

\refitem{Itoh, Y., Langlotz, T., Sutton, J., \& Plopski, A. (2021). Towards Indistinguishable Augmented Reality: A Survey on Optical See-Through Head-Mounted Displays. \textit{ACM Computing Surveys}, 54(6), 1–36. doi:10.1145/3453157}

\refitem{Itti, L., Koch, C., \& Niebur, E. (1998). A model of saliency-based visual attention for rapid scene analysis. \textit{IEEE TPAMI, 20}(11).}

\refitem{Jahn, G., Oehme, A., Krems, J. F., \& Gelau, C. (2005). Peripheral detection as a workload measure in driving: effects of traffic complexity and route guidance system use in a driving study. \textit{Transportation Research Part F, 8}(3), 255–275. \url{https://doi.org/10.1016/j.trf.2005.04.009}}

\refitem{Jocher, G., \& Qiu, J. (2024). \textit{Ultralytics YOLO11} (Version 11.0.0) [Computer software]. Ultralytics. \url{https://github.com/ultralytics/ultralytics}. Citation per Ultralytics' own CITATION.cff / docs (\url{docs.ultralytics.com/models/yolo11/}); DOI not yet assigned by Ultralytics at time of writing.}

\refitem{Kalkofen, D., Mendez, E., \& Schmalstieg, D. (2007). Interactive Focus and Context Visualization for Augmented Reality. \textit{ISMAR 2007}, 191–200. doi:10.1109/ISMAR.2007.4538846}

\refitem{Kim, H. K., Park, J., Choi, Y., \& Choe, M. (2018). Virtual reality sickness questionnaire (VRSQ): motion sickness measurement index in a virtual reality environment. \textit{Applied Ergonomics, 69}, 66–73. \url{https://www.sciencedirect.com/science/article/abs/pii/S000368701730282X}}

\refitem{Laghari, M. K., Shaikh, A. A., Khan, F., \& Siddiqui, A. G. (2025). Native mixed reality compositing on Meta Quest 3: a quantitative feasibility study of ARM-based SoCs and thermal headroom. \textit{arXiv 2509.18929}. \url{https://arxiv.org/abs/2509.18929}}

\refitem{Laghari, M. K., Shaikh, A. A., Khan, F., \& Siddiqui, A. G. (2025). \textit{Native Mixed Reality Compositing on Meta Quest 3: A Quantitative Feasibility Study of ARM-Based SoCs and Thermal Headroom.} arXiv:2509.18929.}

\refitem{Lakens, D. (2017). Equivalence tests: a practical primer for t tests, correlations, and meta-analyses. \textit{Social Psychological and Personality Science, 8}(4), 355–362. \url{https://doi.org/10.1177/1948550617697177}}

\refitem{Lee, J., \& Kim, S. (2025). DiminishAR. \textit{CHI 2025}. doi:10.1145/3706598.3713415 (preprint: arXiv:2403.03875)}

\refitem{Lu, W., Duh, B.-L. H., \& Feiner, S. (2012). Subtle Cueing for Visual Search in Augmented Reality. \textit{2012 IEEE International Symposium on Mixed and Augmented Reality (ISMAR)}, 161–166. doi:10.1109/ISMAR.2012.6402553}

\refitem{Magic Leap (2024). \textit{Global/Segmented Dimmer.} Magic Leap Developer Documentation. \url{developer-docs.magicleap.cloud/docs/guides/features/dimmer-feature/}}

\refitem{McLaughlin, A. C., et al. (2025). Cognitive aid design using diminished reality to support selective attention by reducing distraction. \textit{Human Factors, 67}(9), 937–961. \url{https://journals.sagepub.com/doi/10.1177/00187208251325169}}

\refitem{McNamara, A., Bailey, R., \& Grimm, C. (2008). Improving Search Task Performance Using Subtle Gaze Direction. \textit{APGV 2008}, 51–56.}

\refitem{Mendez, E., Feiner, S., \& Schmalstieg, D. (2010). Focus and Context in Mixed Reality by Modulating First Order Salient Features. \textit{Smart Graphics 2010}, LNCS 6133, 232–243.}

\refitem{Meta (2024a). \textit{Passthrough Styling (colour mapping / LUTs).} \url{developers.meta.com/horizon/documentation/unity/unity-customize-passthrough-color-mapping/}}

\refitem{Meta (2024b). \textit{Passthrough Windows (punch-through).} \url{developers.meta.com/horizon/documentation/unity/unity-customize-passthrough-passthrough-windows/}}

\refitem{Meta (2024c). \textit{MR Motifs: Passthrough Transitioning.} \url{developers.meta.com/horizon/documentation/unity/unity-mrmotifs-passthrough-transitioning/}}

\refitem{Meta (2024d). \textit{Meta Quest build v67 release notes: Theatre View.} Meta Community Forums. \url{communityforums.atmeta.com/blog/AnnouncementsBlog/meta-quest-build-v67-release-notes/1215574}}

\refitem{Meta (2025). \textit{Passthrough Camera Access (PCA).} \url{developers.meta.com/horizon/} (Unity PCA documentation)}

\refitem{Meta (2025a). \textit{How to Take Perfetto Traces with Meta Quest Developer Hub.} \url{developers.meta.com/horizon/documentation/unity/ts-perfettoguide/}}

\refitem{Meta (2025b). \textit{Monitor Performance with OVR Metrics Tool.} \url{developers.meta.com/horizon/documentation/unity/ts-ovrmetricstool/}}

\refitem{Mori, S., Ikeda, S., \& Saito, H. (2017). A survey of diminished reality: techniques for visually concealing, eliminating, and seeing through real objects. \textit{IPSJ T-CVA, 9}(17). \url{https://link.springer.com/article/10.1186/s41074-017-0028-1}}

\refitem{Murph, I., McDonald, M., Richardson, K., Wilkinson, M., Robertson, S., Karunakaran, A., Gandy Coleman, M., Byrne, V., \& McLaughlin, A. C. (2021). Diminishing Reality: Potential Benefits and Risks. \textit{Proceedings of the Human Factors and Ergonomics Society Annual Meeting}, 65(1), 164–168. doi:10.1177/1071181321651103}

\refitem{Murph, I., Richardson, K., \& McLaughlin, A. C. (2022). Methods of Training to Overcome Distraction Via Diminished Reality. \textit{Proceedings of the Human Factors and Ergonomics Society Annual Meeting}, 66(1), 1844–1848. doi:10.1177/1071181322661134}

\refitem{Norouzi, N., Bruder, G., \& Welch, G. (2018). Assessing Vignetting as a Means to Reduce VR Sickness During Amplified Head Rotations. \textit{15th ACM Symposium on Applied Perception (SAP '18).} doi:10.1145/3225153.3225162}

\refitem{Obscuring Undesirable Individuals to Alleviate Social Discomfort Using Diminished Reality. \textit{CHI 2026}. \url{https://dl.acm.org/doi/10.1145/3772318.3790918}}

\refitem{Patents US 12548271 B2 (Apple Inc., granted 2026-02-10) — "Attention control in multi-user environments" (title confirmed verbatim via Google Patents); US 12524072 B2 (Samsung Electronics, granted 2026-01-13) — "Providing a pass-through view of a real-world environment for a virtual reality headset for a user interaction with real world objects", paraphrased here as "blurred external video feed in VR". Both numbers confirmed valid and correctly assigned; claim-level adjacency only, no implementation evidence in either case.}

\refitem{Patney, A., et al. (2016). Towards Foveated Rendering for Gaze-Tracked Virtual Reality. \textit{ACM Transactions on Graphics}, 35(6). doi:10.1145/2980179.2980246}

\refitem{Rusch, M. L., Schall, M. C. Jr., Gavin, P., Lee, J. D., Dawson, J. D., Vecera, S., \& Rizzo, M. (2013). Directing driver attention with augmented reality cues. \textit{Transportation Research Part F: Traffic Psychology and Behaviour, 16}, 127–137. doi:10.1016/j.trf.2012.08.007}

\refitem{Sitzmann, V., Serrano, A., Pavel, A., Agrawala, M., Gutierrez, D., Masia, B., \& Wetzstein, G. (2018). Saliency in VR: How Do People Explore Virtual Environments? \textit{IEEE TVCG}, 24(4), 1633–1642. (arXiv:1612.04335)}

\refitem{Sutton, J., Langlotz, T., Plopski, A., Zollmann, S., Itoh, Y., \& Regenbrecht, H. (2022). Look over there! Investigating Saliency Modulation for Visual Guidance with Augmented Reality Glasses. \textit{UIST 2022}. \url{https://dl.acm.org/doi/10.1145/3526113.3545633}}

\refitem{Syiem, B. V., Kelly, R. M., Goncalves, J., Velloso, E., \& Dingler, T. (2021). Impact of Task on Attentional Tunneling in Handheld Augmented Reality. \textit{CHI 2021}. doi:10.1145/3411764.3445580}

\refitem{Tariq, T., Tursun, C., \& Didyk, P. (2022). Noise-based Enhancement for Foveated Rendering. \textit{ACM Transactions on Graphics.} doi:10.1145/3528223.3530101}

\refitem{Tsandilas, T., \& Casiez, G. (2024). The Illusory Promise of the Aligned Rank Transform. \textit{Journal of Visualization and Interaction} (under review). \url{https://statransform.github.io/jovi/}}

\refitem{Tursun, O. T., et al. (2019). Luminance-Contrast-Aware Foveated Rendering. \textit{ACM Transactions on Graphics}, 38(4). doi:10.1145/3306346.3322985}

\refitem{Unity Technologies (2024). \textit{Unity Sentis} (on-device neural network inference package). \url{https://unity.com/products/sentis}. Note: Unity has since renamed this package "Unity Inference Engine" (\texttt{com.unity.ai.inference}, v2.6 as of 2026); cited here under the "Sentis" name because that is the name under which it was integrated into this dissertation's pipeline.}

\refitem{Veas, E., Mendez, E., Feiner, S., \& Schmalstieg, D. (2011). Directing Attention and Influencing Memory with Visual Saliency Modulation. \textit{CHI 2011}, 1471–1480. doi:10.1145/1978942.1979158}

\refitem{Victor, T. W., Harbluk, J. L., \& Engström, J. A. (2005). Sensitivity of eye-movement measures to in-vehicle task difficulty. \textit{Transportation Research Part F, 8}(2), 167–190. \url{https://doi.org/10.1016/j.trf.2005.04.014}}

\refitem{Waldner, M., Le Muzic, M., Bernhard, M., Purgathofer, W., \& Viola, I. (2014). Attractive Flicker — Guiding Attention in Dynamic Narrative Visualizations. \textit{IEEE TVCG}, 20(12), 2456–2465. doi:10.1109/TVCG.2014.2346352}

\refitem{Wang, G., Gan, Q., \& Li, Y. (2020). Research on Attention-guiding Methods in Cinematic Virtual Reality Based on Eye Tracking Analysis. \textit{2020 International Conference on Innovation Design and Digital Technology (ICIDDT)}. doi:10.1109/ICIDDT52279.2020.00020}

\refitem{Wang, L., Shi, X., \& Liu, Y. (2022). Foveated Rendering: A State-of-the-Art Survey. arXiv:2211.07969.}

\refitem{Wang, J., Ping, S., Xu, K., Li, Y., \& Liang, H.-N. (2026). The perceptual gap between video see-through displays and natural human vision. \textit{arXiv 2601.02805}. \url{https://arxiv.org/pdf/2601.02805}}

\refitem{Wobbrock, J. O., Findlater, L., Gergle, D., \& Higgins, J. J. (2011). The aligned rank transform for nonparametric factorial analyses using only ANOVA procedures. In \textit{Proceedings of CHI 2011} (pp. 143–146). ACM. doi:10.1145/1978942.1978963}

\refitem{Wu, F., \& Suma Rosenberg, E. (2022). Adaptive Field-of-view Restriction: Limiting Optical Flow to Mitigate Cybersickness in Virtual Reality. \textit{VRST 2022}. doi:10.1145/3562939.3565611 (\textbf{corrected 2026-08-05}: the title previously given here — "Asymmetric FOV Restriction Using Optic Flow" — does not match this DOI. That title belongs to a different Wu/Suma Rosenberg paper, \textit{Don't Block the Ground: Reducing Discomfort in Virtual Reality with an Asymmetric Field-of-View Restrictor} (SUI 2021, doi:10.1145/3485279.3485284). The DOI was correct; the title has been fixed to match it.)}

\refitem{Yao, L., DeVincenzi, A., Pereira, A., \& Ishii, H. (2013). FocalSpace: Multimodal Activity Tracking, Synthetic Blur and Adaptive Presentation for Video Conferencing. \textit{ACM SUI 2013}. doi:10.1145/2491367.2491377}
